\chapter{Những sự thật về tiền bạc}
\markboth{Những sự thật về tiền bạc}{Truths About Money}


\section{Không có tiền rất khó sống}
\begin{truyen}{Không có tiền rất khó sống}{Life lesson}
\chuhoa{A}{nh từng không có tiền đóng học, không tiền thuê nhà. Mọi thứ đều khó. Không có tiền rất khó sống — đó là sự thật không cần che giấu.}
\textit{(He had no money for school, no money for rent. Everything was hard. Without money life is very hard — that's a truth we don't need to hide.)}
\end{truyen}

\begin{baihoc}
Cần tiền để sống — nhưng đừng để tiền thành chủ nhân.
\textit{(You need money to live — but don't let money become the master.)}
\end{baihoc}

\begin{ghinhoanh}
\textbf{Short English to remember:}

You need money to live — but don't let money become the master.

\end{ghinhoanh}

\begin{tuhocnhanh}
1. Câu chuyện này gợi cho bạn suy nghĩ gì? \textit{What does this story make you think?}

2. Bạn sẽ nhớ bài học nào nhất? \textit{Which lesson will you remember most?}
\end{tuhocnhanh}
\ngancach


\section{Nhưng có tiền chưa chắc hạnh phúc}
\begin{truyen}{Nhưng có tiền chưa chắc hạnh phúc}{Life lesson}
\chuhoa{M}{ột người giàu có nói: tôi có thể mua nhiều thứ, nhưng tôi không mua được giấc ngủ ngon và sự bình yên. Có tiền chưa chắc hạnh phúc.}
\textit{(A wealthy man said: I can buy many things, but I can't buy good sleep and peace. Having money doesn't guarantee happiness.)}
\end{truyen}

\begin{baihoc}
Tiền giải quyết nhiều chuyện — nhưng không giải quyết hết chuyện lòng.
\textit{(Money solves many things — but not everything in the heart.)}
\end{baihoc}

\begin{ghinhoanh}
\textbf{Short English to remember:}

Money solves many things — but not everything in the heart.

\end{ghinhoanh}

\begin{tuhocnhanh}
1. Câu chuyện này gợi cho bạn suy nghĩ gì? \textit{What does this story make you think?}

2. Bạn sẽ nhớ bài học nào nhất? \textit{Which lesson will you remember most?}
\end{tuhocnhanh}
\ngancach


\section{Người nghèo thường giúp nhau hơn}
\begin{truyen}{Người nghèo thường giúp nhau hơn}{Life lesson}
\chuhoa{Ở}{ khu anh ở, khi nhà nào khó khăn, hàng xóm góp gạo, góp tiền. Người nghèo thường giúp nhau hơn — vì họ hiểu cảm giác thiếu thốn.}
\textit{(In his neighborhood, when a family was in need, neighbors gave rice and money. The poor often help each other more — because they know what it's like to lack.)}
\end{truyen}

\begin{baihoc}
Nghèo không có nghĩa là ích kỷ — đôi khi ngược lại.
\textit{(Being poor doesn't mean being selfish — sometimes the opposite.)}
\end{baihoc}

\begin{ghinhoanh}
\textbf{Short English to remember:}

Being poor doesn't mean being selfish — sometimes the opposite.

\end{ghinhoanh}

\begin{tuhocnhanh}
1. Câu chuyện này gợi cho bạn suy nghĩ gì? \textit{What does this story make you think?}

2. Bạn sẽ nhớ bài học nào nhất? \textit{Which lesson will you remember most?}
\end{tuhocnhanh}
\ngancach


\section{Tiền làm lộ ra tính cách thật}
\begin{truyen}{Tiền làm lộ ra tính cách thật}{Life lesson}
\chuhoa{K}{hi anh vay được tiền, vài người bạn thân biến mất. Khi anh trả xong và có dư, họ quay lại. Tiền làm lộ ra tính cách thật.}
\textit{(When he borrowed money, some close friends disappeared. When he paid back and had extra, they returned. Money reveals true character.)}
\end{truyen}

\begin{baihoc}
Đừng dùng tiền để thử bạn — nhưng khi tiền xuất hiện, hãy mở mắt.
\textit{(Don't use money to test friends — but when money appears, keep your eyes open.)}
\end{baihoc}

\begin{ghinhoanh}
\textbf{Short English to remember:}

Don't use money to test friends — but when money appears, keep yo...

\end{ghinhoanh}

\begin{tuhocnhanh}
1. Câu chuyện này gợi cho bạn suy nghĩ gì? \textit{What does this story make you think?}

2. Bạn sẽ nhớ bài học nào nhất? \textit{Which lesson will you remember most?}
\end{tuhocnhanh}
\ngancach


\section{Tiền không mua được sự tôn trọng thật}
\begin{truyen}{Tiền không mua được sự tôn trọng thật}{Life lesson}
\chuhoa{A}{nh từng nghĩ nhiều tiền sẽ khiến mọi người nể. Sau này anh thấy: sự tôn trọng thật đến từ cách mình đối xử với người, không phải từ ví.}
\textit{(He once thought a lot of money would make everyone respect him. Later he saw: real respect comes from how you treat people, not from your wallet.)}
\end{truyen}

\begin{baihoc}
Tôn trọng mua bằng tiền thường không bền.
\textit{(Respect bought with money usually doesn't last.)}
\end{baihoc}

\begin{ghinhoanh}
\textbf{Short English to remember:}

Respect bought with money usually doesn't last.

\end{ghinhoanh}

\begin{tuhocnhanh}
1. Câu chuyện này gợi cho bạn suy nghĩ gì? \textit{What does this story make you think?}

2. Bạn sẽ nhớ bài học nào nhất? \textit{Which lesson will you remember most?}
\end{tuhocnhanh}
\ngancach


\section{Người giàu cũng có nỗi sợ}
\begin{truyen}{Người giàu cũng có nỗi sợ}{Life lesson}
\chuhoa{M}{ột doanh nhân nói: tôi sợ mất trắng, sợ bệnh tật, sợ con cháu tranh giành. Người giàu cũng có nỗi sợ — chỉ khác loại.}
\textit{(A businessman said: I'm afraid of losing everything, of illness, of my children fighting. The rich have fears too — just different kinds.)}
\end{truyen}

\begin{baihoc}
Ai cũng có nỗi sợ — tiền không xóa được hết.
\textit{(Everyone has fears — money can't erase them all.)}
\end{baihoc}

\begin{ghinhoanh}
\textbf{Short English to remember:}

Everyone has fears — money can't erase them all.

\end{ghinhoanh}

\begin{tuhocnhanh}
1. Câu chuyện này gợi cho bạn suy nghĩ gì? \textit{What does this story make you think?}

2. Bạn sẽ nhớ bài học nào nhất? \textit{Which lesson will you remember most?}
\end{tuhocnhanh}
\ngancach


\section{Tiền đến chậm nhưng đi rất nhanh}
\begin{truyen}{Tiền đến chậm nhưng đi rất nhanh}{Life lesson}
\chuhoa{A}{nh tiết kiệm cả năm mới mua được chiếc xe. Một lần tai nạn, tiền sửa và viện phí làm anh trắng tay. Tiền đến chậm nhưng đi rất nhanh.}
\textit{(He saved a year to buy a car. One accident, repair and hospital bills left him with nothing. Money comes slowly but goes very fast.)}
\end{truyen}

\begin{baihoc}
Tiền kiếm khó — nên dùng có ý thức.
\textit{(Money is hard to earn — so use it with care.)}
\end{baihoc}

\begin{ghinhoanh}
\textbf{Short English to remember:}

Money is hard to earn — so use it with care.

\end{ghinhoanh}

\begin{tuhocnhanh}
1. Câu chuyện này gợi cho bạn suy nghĩ gì? \textit{What does this story make you think?}

2. Bạn sẽ nhớ bài học nào nhất? \textit{Which lesson will you remember most?}
\end{tuhocnhanh}
\ngancach


\section{Kiếm tiền khó hơn tiêu tiền}
\begin{truyen}{Kiếm tiền khó hơn tiêu tiền}{Life lesson}
\chuhoa{A}{nh mất năm năm để dành dụm, nhưng chỉ vài tháng để tiêu hết khi đầu tư sai. Kiếm tiền khó hơn tiêu tiền.}
\textit{(It took him five years to save, but only a few months to spend it all on a bad investment. Earning is harder than spending.)}
\end{truyen}

\begin{baihoc}
Tiêu dễ, kiếm khó — đừng đảo ngược thói quen.
\textit{(Spending is easy, earning is hard — don't reverse the habit.)}
\end{baihoc}

\begin{ghinhoanh}
\textbf{Short English to remember:}

Spending is easy, earning is hard — don't reverse the habit.

\end{ghinhoanh}

\begin{tuhocnhanh}
1. Câu chuyện này gợi cho bạn suy nghĩ gì? \textit{What does this story make you think?}

2. Bạn sẽ nhớ bài học nào nhất? \textit{Which lesson will you remember most?}
\end{tuhocnhanh}
\ngancach


\section{Nợ tiền dễ mất bạn}
\begin{truyen}{Nợ tiền dễ mất bạn}{Life lesson}
\chuhoa{K}{hi anh nợ tiền bạn, bạn ấy xa dần. Khi anh trả xong, tình bạn không còn như trước. Nợ tiền dễ mất bạn.}
\textit{(When he owed his friend money, the friend drifted away. When he paid back, the friendship was never the same. Debt easily loses friends.)}
\end{truyen}

\begin{baihoc}
Nợ tiền không chỉ là nợ số — còn là nợ tình.
\textit{(Debt is not just numbers — it's also relationship.)}
\end{baihoc}

\begin{ghinhoanh}
\textbf{Short English to remember:}

Debt is not just numbers — it's also relationship.

\end{ghinhoanh}

\begin{tuhocnhanh}
1. Câu chuyện này gợi cho bạn suy nghĩ gì? \textit{What does this story make you think?}

2. Bạn sẽ nhớ bài học nào nhất? \textit{Which lesson will you remember most?}
\end{tuhocnhanh}
\ngancach


\section{Tiền là công cụ, không phải mục đích}
\begin{truyen}{Tiền là công cụ, không phải mục đích}{Life lesson}
\chuhoa{A}{nh từng đuổi theo tiền như mục đích. Khi có đủ, anh mới hiểu: tiền là công cụ để sống tốt hơn, không phải điểm đến.}
\textit{(He once chased money as a goal. When he had enough, he understood: money is a tool to live better, not the destination.)}
\end{truyen}

\begin{baihoc}
Tiền là phương tiện — đừng biến nó thành chủ.
\textit{(Money is a means — don't make it the master.)}
\end{baihoc}

\begin{ghinhoanh}
\textbf{Short English to remember:}

Money is a means — don't make it the master.

\end{ghinhoanh}

\begin{tuhocnhanh}
1. Câu chuyện này gợi cho bạn suy nghĩ gì? \textit{What does this story make you think?}

2. Bạn sẽ nhớ bài học nào nhất? \textit{Which lesson will you remember most?}
\end{tuhocnhanh}
\ngancach
